% ============================================================
%  PE5 -- Unidad 5 | PLANTILLA DE TRABAJO DIVIDIDO
%  MundiPets -- Plataforma de Adopción y Cruza Responsable de Mascotas
%  Asignatura: Ingeniería de Requisitos (ISR-401) -- UTEQ 2026-2027 PPA
%  Formato: IEEE onecolumn | Biber + biblatex (estilo IEEE)
%  Preámbulo trasplantado y adaptado del ERS_SRS_2B_v2_0.tex del PFC,
%  para que la PE5 y el ERS compartan exactamente el mismo motor de
%  compilación (mismo compilador, mismo estilo de citas, misma
%  numeración de tablas/figuras). Ver Sección 0.6 y 0.7 del cuerpo.
%
%  COMO USAR ESTE ARCHIVO
%  ----------------------
%  Cada bloque  indica QUE debe producir un
%  integrante en esa sección. El responsable borra la caja 
%  y la reemplaza por el contenido real. Cuando no quede ninguna caja
%  gris en el documento, la PE5 está terminada.
%
%  COMPILACIÓN (requiere referencias.bib en la misma carpeta):
%      pdflatex PE5_U5_Trabajo_Dividido.tex
%      biber    PE5_U5_Trabajo_Dividido
%      pdflatex PE5_U5_Trabajo_Dividido.tex
%      pdflatex PE5_U5_Trabajo_Dividido.tex
% ============================================================
\documentclass[journal,onecolumn]{IEEEtran}

% ── Codificación e idioma ───────────────────────────────────
\usepackage[T1]{fontenc}
\usepackage[utf8]{inputenc}
\usepackage[spanish, es-tabla]{babel}
\usepackage[autostyle=true, spanish=mexican]{csquotes}
\usepackage[htt]{hyphenat}   % permite guionar dentro de \texttt (evita Overfull)

% ── Tablas ──────────────────────────────────────────────────
\usepackage{booktabs}
\usepackage{tabularx}
\usepackage{array}
\usepackage{multirow}
\usepackage{longtable}
\newcolumntype{C}[1]{>{\centering\arraybackslash}p{#1}}
\newcolumntype{L}[1]{>{\raggedright\arraybackslash}p{#1}}
\renewcommand{\thetable}{\arabic{table}}
\def\tablename{Tabla}

% ── Figuras ─────────────────────────────────────────────────
\usepackage{graphicx}
\graphicspath{{figuras/}{./}}
\usepackage{float}
\usepackage{pdflscape}   % matriz de trazabilidad en horizontal

% ── Color ───────────────────────────────────────────────────
\usepackage{xcolor}
\definecolor{uteqverde}{RGB}{0,110,60}
\definecolor{codegray}{gray}{0.95}
\definecolor{codegreen}{rgb}{0.0,0.5,0.0}
\definecolor{codeblue}{rgb}{0.0,0.0,0.6}
\definecolor{codepurple}{rgb}{0.5,0.0,0.35}

% ── Geometría ───────────────────────────────────────────────
\usepackage{geometry}
\geometry{a4paper, top=2.5cm, bottom=2.5cm, left=2.5cm, right=2.5cm,
	headheight=14pt, footskip=1.2cm}

% ── Encabezado y pie ────────────────────────────────────────
\usepackage{fancyhdr}
\pagestyle{fancy}
\fancyhf{}
\renewcommand{\headrulewidth}{0.4pt}
\fancyhead[L]{\small PE5 -- Unidad 5 -- Ingeniería de Requisitos}
\fancyhead[R]{\small MundiPets -- Trabajo dividido}
\fancyfoot[C]{\thepage}

% ── Espaciado y listas ──────────────────────────────────────
\usepackage{parskip}
\setlength{\parskip}{4pt}
\setlength{\emergencystretch}{3em}   % tolerancia extra de línea
\makeatletter
\@ifundefined{labelindent}{}{\let\labelindent\relax}
\makeatother
\usepackage{enumitem}
\usepackage{amssymb}
\usepackage{amsmath}

% ── Títulos de sección ──────────────────────────────────────
\usepackage{titlesec}
\renewcommand{\thesection}{\arabic{section}}
\renewcommand{\thesubsection}{\thesection.\arabic{subsection}}
\renewcommand{\thesubsubsection}{\thesubsection.\arabic{subsubsection}}
\titleformat{\section}{\normalfont\Large\bfseries}{\thesection}{1em}{}
\titleformat{\subsection}{\normalfont\large\bfseries}{\thesubsection}{1em}{}
\titleformat{\subsubsection}{\normalfont\normalsize\bfseries}{\thesubsubsection}{1em}{}
\setcounter{secnumdepth}{4}
\setcounter{tocdepth}{4}

% ── Leyendas ────────────────────────────────────────────────
\usepackage{caption}
\captionsetup{font=small, labelfont=bf}

% ── Bibliografía (biber + estilo IEEE) ──────────────────────
% Mismo motor que el ERS del PFC: ver Sección 0.6 y 0.7 del cuerpo
% para el orden exacto de compilación.
\usepackage[
backend=biber,
style=ieee,
citestyle=numeric-comp,
isbn=true,
doi=true
]{biblatex}
\addbibresource{referencias.bib}

% ── Enlaces ─────────────────────────────────────────────────
\usepackage[hidelinks,colorlinks=false]{hyperref}
\usepackage{url}

% ============================================================
%  CAJA DE GUÍA (fondo gris claro, borde negro)
% ============================================================
\usepackage{mdframed}
\mdfdefinestyle{guia}{
	linecolor=black,
	linewidth=0.6pt,
	backgroundcolor=gray!10,
	innertopmargin=6pt,
	innerbottommargin=6pt,
	innerleftmargin=8pt,
	innerrightmargin=8pt,
	skipabove=4pt,
	skipbelow=4pt
}

% 
\newcommand{\guia}[2]{%
	\begin{mdframed}[style=guia]
		\noindent\textbf{Responsable: #1}\\[2pt]
		\textit{#2}
	\end{mdframed}
}

% -- Atajos de nombres (evita erratas al repetirlos) ---------
\newcommand{\rFUE}{FUERTES ARRAES, Edson Daniel}
\newcommand{\rGUT}{GUTIÉRREZ ORTEGA, Génesis Adriana}
\newcommand{\rMEN}{MENDOZA PÁRRAGA, Andy Johel}
\newcommand{\rMOR}{MORALES SÁNCHEZ, Gary Alejandro}
\newcommand{\rNIE}{NIEVES SÁNCHEZ, Jimmy Samuel}
\newcommand{\rTODOS}{TODO EL EQUIPO}

% ============================================================
\begin{document}
	
	% ============================================================
	% PORTADA (mismo formato que el ERS del PFC, Apéndice de portada)
	% ============================================================

	\section{Requisitos de los Componentes de Inteligencia Artificial}
\label{sec:requisitos_ia}

% -- 7.1 --------------------------------------------------------
\subsection{Identificación y justificación de los componentes}

MundiPets incorpora tres componentes de apoyo basados en inteligencia artificial, ya
declarados en el alcance del sistema y en el requisito de diseño RD-01 (integración de
componentes de inteligencia artificial). Ninguno de los tres sustituye el criterio de un
profesional veterinario ni toma una decisión definitiva por sí solo: los tres producen una
salida orientativa que el usuario o un tercero (veterinaria, moderador) puede revisar. Esta
restricción ya está declarada como límite explícito del sistema desde PE1--PE2 y se mantiene
como marco de todo lo que sigue en esta sección.

\textbf{IA-01 --- Evaluación de compatibilidad para cruza.} Responde a RF-06, cuyo origen
documental combina evidencia de propietarios, interesados en cruza y veterinarias (EV-01 a
EV-08, EV-11, EV-13 a EV-16): los entrevistados señalaron reiteradamente la dificultad de
evaluar de forma informal si dos mascotas son compatibles para un proceso de cruza responsable,
considerando raza, edad, estado de salud, antecedentes genéticos y parentesco. Es, además, el
único de los tres componentes con un derecho legal directamente asociado: el Art.\,20 de la
LOPDP (RL-06) reconoce al titular el derecho a no ser objeto de una decisión basada únicamente
en valoraciones automatizadas sin recibir una explicación motivada, lo que ya obligó en PE3--PE4
a declarar RNF-06 (exactitud) y RNF-13 (explicabilidad) sobre este mismo componente.

\textbf{IA-02 --- Validación de imágenes de perfil.} Responde a RF-17, motivado por evidencia de
propietarios y veterinarias (EV-01, EV-03, EV-05, EV-07, EV-08) sobre la presencia de
fotografías irrelevantes, borrosas o no relacionadas con el propósito de la plataforma en las
publicaciones de mascotas. El componente ya cuenta con RNF-07 (precisión de clasificación) y
RNF-14 (explicabilidad del rechazo) definidos desde PE3--PE4, lo que permite construir su ficha
y sus requisitos de IA sin partir de cero.

\textbf{IA-03 --- Moderación de mensajes.} Responde a RF-07 (sistema de mensajería), sobre cuya
arquitectura ya se especificó en PE3 un componente de moderación de mensajes que clasifica el
contenido antes de su entrega, motivado por la evidencia de coordinación entre propietarios e
interesados (EV-02, EV-07, EV-08). RNF-15 (explicabilidad de la moderación) ya exige que, ante
un mensaje bloqueado, el sistema muestre la categoría de infracción detectada antes de impedir
el envío.

Los tres componentes comparten el mismo límite declarado en el alcance del sistema: ninguno
reemplaza la evaluación veterinaria profesional, y ninguna de sus salidas es definitiva sin
posibilidad de revisión humana. Un solo componente hubiera dejado C3 en N2 según la guía de
esta práctica; los tres cuentan además con evidencia documental propia y no responden a una
moda tecnológica, sino a necesidades identificadas directamente en PE1--PE2 y ya incorporadas
como requisitos del ERS desde PE3.

% -- 7.2 --------------------------------------------------------
\subsection{Ficha del componente IA-01}

\begin{table}[H]
	\caption{Ficha del componente de inteligencia artificial IA-01 --- Evaluación de compatibilidad para cruza.}
	\label{tab:ficha_ia01}
	\small
	\begin{tabularx}{\textwidth}{|L{4.2cm}|X|}
		\hline
		\textbf{Campo} & \textbf{Contenido} \\
		\hline
		Identificador y nombre & IA-01 --- Evaluación de compatibilidad para cruza \\
		\hline
		Tarea y tipo & Recomendación (\textit{scoring} de compatibilidad con explicación causal) \\
		\hline
		Entradas y salidas & \textbf{Entrada:} raza, edad, sexo, estado de salud, antecedentes genéticos, parentesco e historial médico de las dos mascotas seleccionadas. \textbf{Salida:} nivel de compatibilidad (viable / riesgoso / requiere revisión veterinaria) junto con una explicación causal de máximo 60 palabras (RNF-13) y la advertencia de que el resultado es orientativo. \\
		\hline
		Consumidor del resultado & Propietario de mascota e interesado en cruza, para decidir si continúan o no con el proceso de cruza (RF-06); no autoriza ni rechaza el proceso de forma definitiva. \\
		\hline
		Datos de entrenamiento & \textbf{Origen:} historiales clínicos y genéticos registrados en la plataforma más evaluaciones veterinarias etiquetadas durante la fase de pruebas. \textbf{Volumen estimado:} conjunto inicial de al menos 300 pares de mascotas evaluados por veterinarios para el conjunto de referencia de RNF-06. \textbf{Etiquetado:} veterinario colegiado, sobre las mismas variables declaradas en RF-06. \textbf{Calidad mínima:} historiales con antecedentes genéticos y médicos completos; se descartan registros con campos obligatorios vacíos. \textbf{Sesgos conocidos:} sobrerrepresentación esperable de razas comunes en la zona de despliegue inicial frente a razas poco frecuentes, lo que puede degradar la exactitud en estas últimas. \textbf{Base legal:} consentimiento del titular y protección de datos desde el diseño (Art.\,7, 8 y 39 LOPDP) \cite{lopdp2021,lopdp_reglamento2023}. \textbf{Conservación:} mientras la mascota o el propietario mantengan una cuenta activa en la plataforma. \\
		\hline
		Métricas de éxito & \textbf{Principal:} coincidencia con la evaluación veterinaria $\geq$ 85\,\% (RNF-06). \textbf{Secundarias:} tiempo de generación de la explicación $\leq$ 2\,s y calificación media de comprensión $\geq$ 4/5 en escala Likert (RNF-13). \textbf{Conjunto de evaluación:} muestra representativa de pares de mascotas con evaluación veterinaria de referencia, distinta del conjunto de entrenamiento. \\
		\hline
		Equidad & Métrica: diferencia de tasa de coincidencia con la evaluación veterinaria entre razas comunes y razas poco representadas en los datos. Grupos comparados: razas con más de 30 casos en el conjunto de entrenamiento frente a razas con menos de 30 casos. Brecha máxima tolerada: 10 puntos porcentuales (RNF-IA-04, Sección~\ref{sec:req_ia_tabla}). \\
		\hline
		Explicabilidad & Se explica el resultado de compatibilidad al propietario y al interesado en cruza, en el momento de recibir el resultado, mediante una explicación causal por factores (raza, edad, estado de salud, antecedentes genéticos) de máximo 60 palabras, conforme a RNF-13 y al derecho a explicación de decisiones automatizadas de RL-06. \\
		\hline
		Riesgos y \textit{fallback} & Si el modelo no está disponible o no alcanza el umbral mínimo de confianza en un caso concreto, el sistema muestra el mensaje ``compatibilidad no determinada, se recomienda evaluación veterinaria'' en lugar de forzar una recomendación; el proceso de cruza puede continuar manualmente sin el componente. Una recomendación errónea no bloquea la decisión final del usuario, que conserva la responsabilidad última. \\
		\hline
		Clasificación de riesgo & Riesgo limitado conforme al enfoque basado en riesgo del Reglamento (UE) 2024/1689 \cite{euaiact2024}: no es un sistema de identificación biométrica ni de puntuación social, pero al influir en una decisión con efecto sobre el bienestar animal exige las obligaciones de transparencia y explicabilidad ya declaradas en RNF-13 y RL-06. \\
		\hline
		Plan de monitoreo & Ver Sección~\ref{sec:monitoreo_ia}: indicador principal de coincidencia con evaluación veterinaria, medido mensualmente sobre una muestra de recomendaciones auditadas. \\
		\hline
		Trazas & RF-06, RNF-06, RNF-13, RD-01, RL-06. \\
		\hline
	\end{tabularx}
\end{table}

\begin{table}[H]
	\caption{Ficha del componente de inteligencia artificial IA-02 --- Validación de imágenes de perfil.}
	\label{tab:ficha_ia02}
	\small
	\begin{tabularx}{\textwidth}{|L{4.2cm}|X|}
		\hline
		\textbf{Campo} & \textbf{Contenido} \\
		\hline
		Identificador y nombre & IA-02 --- Validación de imágenes de perfil \\
		\hline
		Tarea y tipo & Clasificación (visión: imagen válida / inválida, con motivo de rechazo) \\
		\hline
		Entradas y salidas & \textbf{Entrada:} imagen cargada por el usuario al perfil de una mascota. \textbf{Salida:} resultado binario (aceptada / rechazada) y, si es rechazada, el motivo específico (contenido inapropiado, imagen borrosa, tipo de fotografía incorrecto) en máximo 30 palabras (RNF-14). \\
		\hline
		Consumidor del resultado & Propietario de mascota, que recibe el resultado al momento de publicar; veterinaria, en el caso de imágenes asociadas a validación de información médica. \\
		\hline
		Datos de entrenamiento & \textbf{Origen:} conjunto etiquetado de fotografías de mascotas (válidas) y de imágenes irrelevantes, ofensivas o de baja calidad (inválidas), construido para las pruebas de RNF-07. \textbf{Volumen estimado:} conjunto de evaluación con representación equilibrada de ambas clases, suficiente para sostener la meta de RNF-07. \textbf{Etiquetado:} manual, por el equipo responsable de moderación de contenido, siguiendo la política de contenido de la plataforma. \textbf{Calidad mínima:} imágenes con resolución legible; se descartan archivos corruptos o no soportados antes del etiquetado. \textbf{Sesgos conocidos:} posible menor exactitud en mascotas de especies poco representadas en el conjunto inicial (por ejemplo, especies distintas a perros y gatos). \textbf{Base legal:} consentimiento del titular al publicar contenido e interés legítimo de la plataforma en mantener la política de contenido (Art.\,7 y 8 LOPDP) \cite{lopdp2021}. \textbf{Conservación:} mientras la publicación permanezca activa en la plataforma. \\
		\hline
		Métricas de éxito & \textbf{Principal:} clasificación correcta $\geq$ 95\,\% (RNF-07). \textbf{Secundaria:} 100\,\% de los rechazos con motivo específico mostrado y calificación media de comprensión $\geq$ 4/5 (RNF-14). \textbf{Conjunto de evaluación:} lote de imágenes etiquetadas manualmente, distinto del usado para ajustar el clasificador. \\
		\hline
		Equidad & Métrica: diferencia de tasa de rechazo correcto entre especies con mayor representación en el conjunto de entrenamiento (perros, gatos) y especies con menor representación. Grupos comparados: especies con más de 50 imágenes de referencia frente a especies con menos de 50. Brecha máxima tolerada: 10 puntos porcentuales (RNF-IA-05, Sección~\ref{sec:req_ia_tabla}). \\
		\hline
		Explicabilidad & Se explica el motivo del rechazo al propietario, en el momento de intentar publicar la imagen, mediante una explicación por contraste (``se rechazó por X en lugar de Y'') de máximo 30 palabras, conforme a RNF-14. \\
		\hline
		Riesgos y \textit{fallback} & Si el componente no está disponible, el sistema permite la publicación en estado ``pendiente de revisión'' y la deriva a moderación manual, en vez de bloquear indefinidamente al usuario. Un falso rechazo puede apelarse mediante revisión manual; un falso positivo (imagen inapropiada no detectada) queda sujeto al reporte posterior de otros usuarios. \\
		\hline
		Clasificación de riesgo & Riesgo mínimo/limitado conforme al Reglamento (UE) 2024/1689 \cite{euaiact2024}: es un sistema de moderación de contenido, no de identificación de personas; su principal obligación derivada es la transparencia del motivo de rechazo, ya cubierta por RNF-14. \\
		\hline
		Plan de monitoreo & Ver Sección~\ref{sec:monitoreo_ia}: indicador de tasa de clasificación correcta, medido sobre una muestra semanal de imágenes revisadas manualmente como control de calidad. \\
		\hline
		Trazas & RF-17, RNF-07, RNF-14, RD-01. \\
		\hline
	\end{tabularx}
\end{table}

\begin{table}[H]
	\caption{Ficha del componente de inteligencia artificial IA-03 --- Moderación de mensajes.}
	\label{tab:ficha_ia03}
	\small
	\begin{tabularx}{\textwidth}{|L{4.2cm}|X|}
		\hline
		\textbf{Campo} & \textbf{Contenido} \\
		\hline
		Identificador y nombre & IA-03 --- Moderación de mensajes \\
		\hline
		Tarea y tipo & Clasificación (texto: lenguaje ofensivo / spam / solicitud de pago externo / conversación dentro de tema) \\
		\hline
		Entradas y salidas & \textbf{Entrada:} texto del mensaje enviado entre dos usuarios en el sistema de mensajería (RF-07). \textbf{Salida:} decisión de entrega o bloqueo del mensaje y, si se bloquea, la categoría de infracción detectada (RNF-15) mostrada antes del bloqueo. \\
		\hline
		Consumidor del resultado & El usuario que redacta el mensaje, que recibe la notificación de bloqueo con la categoría antes de que el mensaje se envíe. \\
		\hline
		Datos de entrenamiento & \textbf{Origen:} conjunto etiquetado de mensajes de ejemplo con y sin contenido inapropiado, construido para las pruebas del componente. \textbf{Volumen estimado:} conjunto de evaluación balanceado entre las categorías declaradas (ofensivo, spam, pago externo, fuera de tema, apropiado). \textbf{Etiquetado:} manual, por el equipo de moderación, según la política de contenido de mensajería. \textbf{Calidad mínima:} mensajes en español, del dominio de adopción/cruza de mascotas, sin datos personales de terceros no involucrados. \textbf{Sesgos conocidos:} riesgo de sobre-bloqueo de mensajes con vocabulario coloquial del dominio veterinario que el modelo no ha visto en entrenamiento (por ejemplo, términos médicos poco comunes mal clasificados como spam). \textbf{Base legal:} interés legítimo de la plataforma en prevenir fraude y contenido ofensivo, y consentimiento del titular al aceptar los términos de uso de la mensajería (Art.\,7, 8 y 10 lit.\,e LOPDP) \cite{lopdp2021}. \textbf{Conservación:} la conversación se conserva mientras ambos usuarios mantengan cuenta activa; los mensajes bloqueados no se entregan pero se registran para auditoría de moderación. \\
		\hline
		Métricas de éxito & \textbf{Principal:} 100\,\% de los mensajes moderados con categoría de infracción mostrada antes del bloqueo (RNF-15). \textbf{Secundaria:} calificación media de comprensión de la notificación $\geq$ 4/5 (RNF-15). \textbf{Conjunto de evaluación:} lote de mensajes de prueba etiquetados manualmente, con casos de las cuatro categorías de infracción y de mensajes apropiados. \\
		\hline
		Equidad & Métrica: diferencia de tasa de falsos bloqueos (mensajes apropiados marcados como infracción) entre mensajes que usan vocabulario técnico veterinario y mensajes de lenguaje general. Grupos comparados: mensajes con terminología clínica frente a mensajes sin ella. Brecha máxima tolerada: 10 puntos porcentuales (RNF-IA-04\,bis, Sección~\ref{sec:req_ia_tabla}). \\
		\hline
		Explicabilidad & Se explica al remitente, antes de bloquear el envío, la categoría de infracción detectada mediante una explicación por ejemplos (``mensajes como este suelen clasificarse como X''), conforme a RNF-15. \\
		\hline
		Riesgos y \textit{fallback} & Si el componente no está disponible, los mensajes se entregan sin moderación automática y quedan sujetos al canal de reporte manual entre usuarios, en vez de bloquear por completo la mensajería. Un falso bloqueo puede apelarse solicitando revisión manual del mensaje retenido. \\
		\hline
		Clasificación de riesgo & Riesgo limitado conforme al Reglamento (UE) 2024/1689 \cite{euaiact2024}: modera comunicación entre usuarios sin tomar decisiones sobre acceso a servicios esenciales; su obligación principal es la transparencia de la categoría de bloqueo, ya cubierta por RNF-15. \\
		\hline
		Plan de monitoreo & Ver Sección~\ref{sec:monitoreo_ia}: indicador de tasa de falsos bloqueos, medido mensualmente sobre una muestra de mensajes bloqueados revisados manualmente. \\
		\hline
		Trazas & RF-07, RNF-15, RD-01. \\
		\hline
	\end{tabularx}
\end{table}
\newpage

% -- 7.3 --------------------------------------------------------
\subsection{Requisitos funcionales y no funcionales de IA}
\label{sec:req_ia_tabla}

\begin{longtable}{|L{2.4cm}|c|L{7.2cm}|L{3.6cm}|}
	\caption{Requisitos funcionales y no funcionales de los componentes de IA.}
	\label{tab:req_ia} \\
	\hline
	\textbf{ID} & \textbf{Tipo} & \textbf{Enunciado con umbral y unidad} & \textbf{Método de comprobación} \\
	\hline
	\endfirsthead
	
	\multicolumn{4}{l}{\small\textit{Tabla~\ref{tab:req_ia} (continuación).}} \\
	\hline
	\textbf{ID} & \textbf{Tipo} & \textbf{Enunciado con umbral y unidad} & \textbf{Método de comprobación} \\
	\hline
	\endhead
	
	\hline
	\multicolumn{4}{r}{\small\textit{Continúa en la página siguiente.}} \\
	\endfoot
	
	\hline
	\endlastfoot
	
	\multicolumn{4}{|l|}{\textit{IA-01 --- Evaluación de compatibilidad para cruza}} \\
	\hline
	RF-IA-01 & RF & El sistema deberá generar un resultado de compatibilidad (viable / riesgoso / requiere revisión veterinaria) para todo par de mascotas seleccionado, en un tiempo máximo de 3\,s. & Seleccionar 20 pares de mascotas y medir el tiempo hasta obtener el resultado. \\
	\hline
	RF-IA-02 & RF & El sistema deberá mostrar, junto al resultado de compatibilidad, una explicación causal de máximo 60 palabras que liste los factores considerados (raza, edad, estado de salud, antecedentes genéticos). & Verificar en 20 casos de prueba que la explicación no supere 60 palabras y mencione al menos dos factores. \\
	\hline
	RF-IA-03 & RF & El sistema deberá incluir la advertencia ``resultado orientativo, no sustituye evaluación veterinaria'' en el 100\,\% de los resultados de compatibilidad emitidos. & Revisar 20 resultados emitidos y comprobar la presencia literal de la advertencia. \\
	\hline
	RNF-IA-01 & Rendimiento & El sistema deberá generar recomendaciones de compatibilidad con una coincidencia $\geq$ 85\,\% respecto a la evaluación veterinaria de referencia (RNF-06). & Comparar 100 recomendaciones contra la evaluación de un veterinario sobre el conjunto de prueba. \\
	\hline
	RNF-IA-02 & Rendimiento & El sistema deberá generar la explicación causal en un tiempo no mayor a 2\,s (RNF-13). & Medir el tiempo de generación en 30 solicitudes consecutivas. \\
	\hline
	RNF-IA-03 & Rendimiento & El sistema deberá mantener el componente de compatibilidad disponible $\geq$ 99\,\% del tiempo mensual, mostrando un mensaje de no disponibilidad ante fallas en lugar de bloquear el flujo de cruza. & Registrar el tiempo de actividad del servicio de IA durante un mes de monitoreo. \\
	\hline
	RNF-IA-04 & Equidad & El sistema deberá mantener una diferencia de tasa de coincidencia con la evaluación veterinaria no mayor a 10 puntos porcentuales entre razas con $>$30 casos y razas con $<$30 casos en el conjunto de entrenamiento. & Calcular la tasa de coincidencia por grupo de razas sobre el conjunto de evaluación y comparar la brecha. \\
	\hline
	RNF-IA-05 & Equidad & El sistema deberá mantener una diferencia de tasa de coincidencia no mayor a 10 puntos porcentuales entre mascotas de razas puras y mestizas. & Calcular la tasa de coincidencia por subgrupo (pura / mestiza) sobre el conjunto de evaluación y comparar la brecha. \\
	\hline
	RNF-IA-06 & Explicabilidad & El sistema deberá lograr una calificación media de comprensión de la explicación $\geq$ 4/5 en escala Likert 1--5 (RNF-13). & Aplicar encuesta de comprensión a usuarios técnicos y no técnicos tras mostrar el resultado. \\
	\hline
	\multicolumn{4}{|l|}{\textit{IA-02 --- Validación de imágenes de perfil}} \\
	\hline
	RF-IA-07 & RF & El sistema deberá clasificar toda imagen cargada como aceptada o rechazada en un tiempo máximo de 2\,s. & Cargar 30 imágenes y medir el tiempo de respuesta de la clasificación. \\
	\hline
	RF-IA-08 & RF & El sistema deberá mostrar, cuando una imagen sea rechazada, el motivo específico del rechazo en máximo 30 palabras (RNF-14). & Cargar 15 imágenes inválidas y verificar la presencia y extensión del motivo mostrado. \\
	\hline
	RF-IA-09 & RF & El sistema deberá permitir cargar una nueva imagen de forma inmediata tras un rechazo, sin reiniciar el flujo de publicación. & Rechazar una imagen y comprobar que el usuario puede cargar una nueva sin salir del formulario de publicación. \\
	\hline
	RNF-IA-07 & Rendimiento & El sistema deberá clasificar correctamente $\geq$ 95\,\% de las imágenes del conjunto de evaluación (RNF-07). & Evaluar el componente contra un conjunto de imágenes etiquetadas manualmente y calcular la tasa de acierto. \\
	\hline
	RNF-IA-08 & Rendimiento & El sistema deberá procesar imágenes de hasta 10\,MB sin superar el tiempo de clasificación declarado en RF-IA-07. & Cargar imágenes de distintos tamaños (hasta 10\,MB) y medir el tiempo de clasificación. \\
	\hline
	RNF-IA-09 & Rendimiento & El sistema deberá mantener el componente de validación disponible $\geq$ 99\,\% del tiempo mensual, derivando la imagen a revisión manual en caso de indisponibilidad en lugar de bloquear la publicación. & Registrar el tiempo de actividad del servicio de validación de imágenes durante un mes de monitoreo. \\
	\hline
	RNF-IA-10 & Equidad & El sistema deberá mantener una diferencia de tasa de clasificación correcta no mayor a 10 puntos porcentuales entre especies con $>$50 imágenes de referencia y especies con $<$50. & Calcular la tasa de acierto por especie sobre el conjunto de evaluación y comparar la brecha. \\
	\hline
	RNF-IA-11 & Equidad & El sistema deberá mantener una diferencia de tasa de falso rechazo no mayor a 10 puntos porcentuales entre imágenes con buena y baja iluminación. & Evaluar el componente sobre subconjuntos de imágenes agrupados por condición de iluminación. \\
	\hline
	RNF-IA-12 & Explicabilidad & El sistema deberá mostrar un motivo específico en el 100\,\% de los rechazos, con calificación media de comprensión $\geq$ 4/5 en escala Likert. & Revisar 15 rechazos y encuestar la comprensión del motivo mostrado a los usuarios. \\
	\hline
	\multicolumn{4}{|l|}{\textit{IA-03 --- Moderación de mensajes}} \\
	\hline
	RF-IA-10 & RF & El sistema deberá clasificar todo mensaje enviado como apropiado o infractor (ofensivo, spam, pago externo, fuera de tema) antes de su entrega. & Enviar 30 mensajes de las cuatro categorías de infracción y uno apropiado, y verificar la clasificación. \\
	\hline
	RF-IA-11 & RF & El sistema deberá mostrar, cuando un mensaje sea bloqueado, la categoría de infracción detectada antes de impedir el envío (RNF-15). & Enviar 15 mensajes con contenido inapropiado y verificar que se muestra la categoría antes del bloqueo. \\
	\hline
	RF-IA-12 & RF & El sistema deberá permitir editar y reenviar un mensaje bloqueado sin perder el contenido no infractor del mensaje original. & Bloquear un mensaje, editarlo y comprobar que el sistema conserva el texto no infractor al reenviarlo. \\
	\hline
	RNF-IA-13 & Rendimiento & El sistema deberá emitir la decisión de moderación en un tiempo máximo de 1\,s desde el envío del mensaje. & Enviar 30 mensajes y medir el tiempo entre el envío y la decisión de moderación. \\
	\hline
	RNF-IA-14 & Rendimiento & El sistema deberá mostrar la categoría de infracción en el 100\,\% de los mensajes moderados, antes del bloqueo (RNF-15). & Revisar 15 mensajes bloqueados y comprobar la presencia de la categoría mostrada. \\
	\hline
	RNF-IA-15 & Rendimiento & El sistema deberá mantener el componente de moderación disponible $\geq$ 99\,\% del tiempo mensual, entregando los mensajes sin moderación automática en caso de indisponibilidad en lugar de bloquear la mensajería. & Registrar el tiempo de actividad del servicio de moderación durante un mes de monitoreo. \\
	\hline
	RNF-IA-16 & Equidad & El sistema deberá mantener una diferencia de tasa de falso bloqueo no mayor a 10 puntos porcentuales entre mensajes con vocabulario técnico veterinario y mensajes de lenguaje general. & Evaluar el componente sobre subconjuntos de mensajes agrupados por tipo de vocabulario y comparar la brecha. \\
	\hline
	RNF-IA-17 & Equidad & El sistema deberá mantener una diferencia de tasa de falso bloqueo no mayor a 10 puntos porcentuales entre mensajes largos ($>$200 caracteres) y mensajes cortos. & Evaluar el componente sobre subconjuntos de mensajes agrupados por longitud y comparar la brecha. \\
	\hline
	RNF-IA-18 & Explicabilidad & El sistema deberá lograr una calificación media de comprensión de la notificación de bloqueo $\geq$ 4/5 en escala Likert 1--5. & Aplicar encuesta de comprensión a usuarios tras recibir una notificación de bloqueo. \\
	\hline
\end{longtable}

% -- 7.4 --------------------------------------------------------
\subsection{Ética, privacidad y clasificación de riesgo}
\label{sec:etica_ia}

\textbf{Finalidad.} Cada componente de IA de MundiPets tiene una finalidad determinada y
declarada en su ficha (Tablas~\ref{tab:ficha_ia01}--\ref{tab:ficha_ia03}): apoyar la decisión
de un proceso de cruza (IA-01), mantener la calidad del contenido publicado (IA-02) y proteger
la comunicación entre usuarios (IA-03). Ninguno de los tres reutiliza los datos tratados para
un fin distinto del declarado, conforme al principio de finalidad de la LOPDP \cite{lopdp2021}.

\textbf{Minimización de datos.} IA-01 procesa únicamente las variables clínicas y genéticas
necesarias para evaluar compatibilidad (Tabla~\ref{tab:ficha_ia01}), sin acceder a datos de
contacto o ubicación del propietario. IA-02 procesa la imagen cargada, sin metadatos de
geolocalización asociados. IA-03 procesa el texto del mensaje, sin acceder al historial médico
o genético de las mascotas de los usuarios que conversan. Este alcance limitado de variables
por componente responde directamente al principio de minimización del Art.\,10 LOPDP.

\textbf{Consentimiento.} El tratamiento de datos por los tres componentes queda cubierto por el
consentimiento otorgado al registrarse y al aceptar los términos de uso de la plataforma
(Art.\,7 y 8 LOPDP), en línea con la base legal ya declarada en RD-04 (protección de datos
personales) para el resto del sistema. El usuario puede rechazar el uso de IA-01 y proceder de
forma manual sin perder acceso al resto de funcionalidades del sistema, conforme al
\textit{fallback} declarado en su ficha.

\textbf{Seudonimización.} Los conjuntos de evaluación usados para medir RNF-IA-04, RNF-IA-05,
RNF-IA-10, RNF-IA-11, RNF-IA-16 y RNF-IA-17 (Tabla~\ref{tab:req_ia}) se construyen con
identificadores internos desvinculados del nombre y contacto del propietario, conforme al
mecanismo de seudonimización que desarrolla el Reglamento General de la LOPDP
\cite{lopdp_reglamento2023}; la reidentificación solo es posible mediante el sistema de
producción, no mediante el conjunto de evaluación.

\textbf{Supervisión humana.} Los tres componentes son de apoyo y no de decisión definitiva: la
salida de IA-01 no autoriza ni rechaza un proceso de cruza, la de IA-02 puede derivarse a
revisión manual y la de IA-03 puede apelarse. En los tres casos existe una vía de revisión
humana declarada en el campo ``Riesgos y \textit{fallback}'' de la ficha correspondiente, en
línea con el enfoque de supervisión humana significativa del Reglamento (UE) 2024/1689
\cite{euaiact2024}.

\textbf{Clasificación de riesgo.} Conforme al enfoque basado en niveles de riesgo del
Reglamento (UE) 2024/1689, ninguno de los tres componentes de MundiPets pertenece a la
categoría de riesgo inaceptable ni de alto riesgo definidas en dicho reglamento: no realizan
identificación biométrica, puntuación social ni evaluación de acceso a servicios esenciales.
IA-01 se clasifica como riesgo limitado por influir en una decisión con efecto sobre el
bienestar animal y estar sujeto a un derecho legal de explicación (RL-06); IA-02 e IA-03 se
clasifican como riesgo mínimo/limitado, con la obligación de transparencia ya cubierta por
RNF-14 y RNF-15 respectivamente. Estas clasificaciones se detallan por componente en el campo
``Clasificación de riesgo'' de cada ficha (Tablas~\ref{tab:ficha_ia01}--\ref{tab:ficha_ia03}).

\textbf{Base legal específica de Ecuador.} Al tratarse de datos personales de propietarios y,
en el caso de IA-01, de datos clínicos y genéticos de las mascotas asociadas a una persona
identificable, el tratamiento de los tres componentes se rige por la Ley Orgánica de
Protección de Datos Personales \cite{lopdp2021} y su Reglamento General
\cite{lopdp_reglamento2023}, además del Art.\,20 LOPDP (RL-06) ya declarado como base legal
directa del bloque de explicabilidad de IA-01. Esta base legal no es una declaración genérica:
tiene fuente normativa identificable (artículo y norma), criterio de aceptación (los campos de
las fichas de esta sección) y traza hacia los requisitos correspondientes, igual que cualquier
otro requisito legal del sistema.

% -- 7.5 --------------------------------------------------------
\subsection{Plan de monitoreo y trazado de los requisitos de IA}
\label{sec:monitoreo_ia}

\begin{table}[H]
	\caption{Plan de monitoreo en producción de los componentes de inteligencia artificial.}
	\label{tab:monitoreo_ia}
	\small
	\begin{tabularx}{\textwidth}{|L{1.6cm}|X|L{3.0cm}|L{3.6cm}|}
		\hline
		\textbf{Comp.} & \textbf{Indicador y frecuencia} & \textbf{Umbral de alerta por deriva} & \textbf{Criterio de reentrenamiento o retirada} \\
		\hline
		IA-01 & Coincidencia con evaluación veterinaria (RNF-IA-01), medida mensualmente sobre una muestra auditada de recomendaciones emitidas. & Caída sostenida por debajo de 80\,\% durante dos meses consecutivos. & Reentrenar con el conjunto de casos auditados del período; si tras el reentrenamiento no se recupera el umbral de RNF-06 en el siguiente ciclo, retirar el componente y volver al flujo manual declarado en su \textit{fallback}. \\
		\hline
		IA-02 & Tasa de clasificación correcta (RNF-IA-07), medida semanalmente sobre una muestra de imágenes revisadas manualmente como control de calidad. & Caída sostenida por debajo de 90\,\% durante tres semanas consecutivas. & Reentrenar incorporando los casos mal clasificados detectados en el control de calidad; si la tasa no se recupera en dos ciclos de reentrenamiento, derivar temporalmente todas las imágenes a revisión manual. \\
		\hline
		IA-03 & Tasa de falsos bloqueos (RNF-IA-16, RNF-IA-17), medida mensualmente sobre una muestra de mensajes bloqueados revisados manualmente. & Tasa de falsos bloqueos superior a 15\,\% durante un mes de medición. & Reentrenar con los mensajes mal clasificados identificados en la revisión; si la tasa no mejora, desactivar temporalmente la moderación automática y operar con el canal de reporte manual entre usuarios. \\
		\hline
	\end{tabularx}
\end{table}

Los tres planes de monitoreo comparten un mismo principio: ninguna deriva detectada retira un
componente de forma automática. El umbral de alerta dispara una revisión y, si corresponde, un
reentrenamiento; solo la falta de recuperación tras el reentrenamiento activa el \textit{fallback}
manual ya declarado en la ficha de cada componente (Tablas~\ref{tab:ficha_ia01}--\ref{tab:ficha_ia03}),
de modo que el sistema nunca queda sin una vía de operación cuando un componente de IA se
degrada o deja de estar disponible.

Finalmente, los 27 requisitos de la Tabla~\ref{tab:req_ia} se trazan hacia los RF y RNF
preexistentes del sistema del siguiente modo: los requisitos de IA-01 (RF-IA-01 a RF-IA-03,
RNF-IA-01 a RNF-IA-06) trazan hacia RF-06, RNF-06, RNF-13 y RL-06; los de IA-02 (RF-IA-07 a
RF-IA-09, RNF-IA-07 a RNF-IA-12) trazan hacia RF-17, RNF-07 y RNF-14; y los de IA-03 (RF-IA-10 a
RF-IA-12, RNF-IA-13 a RNF-IA-18) trazan hacia RF-07 y RNF-15. Estas filas se entregan a
MORALES SÁNCHEZ, Gary Alejandro para su integración en la matriz de trazabilidad final
(Sección 6, Matriz de trazabilidad final): ningún requisito de IA queda aislado, porque todos sirven a
una funcionalidad ya especificada o modifican su criterio de aceptación.

	
	\newpage

\end{document}